% !TeX root = ../notes.tex
% Proposed prose version aligned to 07-execute-plan.tex.

\chapter{Generating high-quality chemical data}

\begin{notebox}
  \keypoint{What do you need to learn from this chapter?}
  Modelling and experiment selection depend on experiments that are executed reproducibly, measured appropriately, and recorded with enough context to be reused.
  High throughput is valuable only when it increases relevant information rather than the rate at which unrecognised error is produced.
\end{notebox}

\section{More experiments are not necessarily more information}

The number of experiments is an incomplete measure of a dataset.
A thousand near-identical endpoints may contain less useful information than a smaller set designed to separate competing chemical explanations.
Chemical diversity, measurement density, precision, and the distribution of successful and unsuccessful outcomes all affect information content.

Endpoint measurements answer whether a selected outcome was present after a chosen time.
Reaction-progress measurements also show induction periods, catalyst deactivation, product inhibition, and changes in rate.
The appropriate measurement therefore depends on whether the objective is screening, prediction, optimisation, or mechanistic understanding.

\begin{warningbox}
  \keypoint{Automation amplifies the experimental design.}
  It can amplify informative measurements, but it can also amplify systematic dosing, sampling, or analytical errors.
\end{warningbox}

\section{Practical FAIR chemical data}

The FAIR principles describe data that are findable, accessible, interoperable, and reusable.\mcite{Wilkinson2016}
In chemical experiments, reusability requires more than a table of yields.
The record should identify materials and structures, quantities and units, conditions and timing, vessel or reactor configuration, instrument methods, calibration, processing, and quality-control decisions.

FAIR does not necessarily mean open.
Sensitive or proprietary data can remain access controlled while still using persistent identifiers, shared schemas, defined units, and traceable processing.
Accessibility means that authorised users and machines can retrieve the data and its metadata through a clear procedure.

Unsuccessful experiments are part of the scientific record.
A zero signal, a failed reaction, an instrument fault, and an unprocessed sample are different outcomes and should have different status codes.
Deleting them or merging them into one numerical zero destroys information needed for modelling and reproducibility.

\section{High-throughput and parallel experimentation}

Parallel experimentation uses arrays of reaction vessels or wells to perform many experiments under coordinated conditions.
Liquid handlers can dispense small volumes reproducibly, while miniaturisation reduces reagent use and increases the number of conditions that can be tested.
The same small volumes make evaporation, adsorption, mixing, and dosing errors more important.

A plate layout is part of the experimental design.
Randomisation, controls, blanks, standards, and replicates help separate chemical trends from spatial or temporal artefacts.
Reagent identity and concentration should be checked before a campaign because automation will execute an incorrect stock solution faithfully.

\section{Measurement is part of the experiment}

An automated reaction is not complete when the scheduled heating or mixing ends.
Sampling, quenching, dilution, transfer, analytical measurement, calibration, processing, and quality control determine the recorded response.
Each step can change the sample or introduce uncertainty.

Chromatography can separate products and internal standards, mass spectrometry can provide molecular evidence, NMR can measure composition and structure, and optical spectroscopy can provide rapid time-resolved signals.
No analytical method answers every chemical question.
Orthogonal methods are especially valuable when an automated decision depends on distinguishing conversion, identity, purity, topology, or function.

The status ``no product'' should be used only when the analysis was capable of detecting the expected product and passed its quality checks.
A saturated detector, missing injection, failed peak integration, or ambiguous mass signal is an analytical status rather than a chemical negative result.

\section{Flow as a data-generation platform}

Flow reactors vary conditions continuously and are therefore well suited to systematic data generation.
For an ideal reactor volume $V$ operated at volumetric flow rate $Q$, the nominal residence time is
\begin{equation}
  \tau=\frac{V}{Q}.
\end{equation}
The equation is useful only when the actual mixing, dispersion, phases, and reactor volume are understood.

A condition change at the inlet does not appear instantaneously at the outlet.
Transport lag and dispersion create transitions between experiments, and a sample collected before steady state may contain material produced under several conditions.
Inline or at-line analysis must therefore be synchronised with the physical movement of reaction segments.

Flow should be selected because its control, heat and mass transfer, pressure range, or analytical integration suits the chemical question.
It is not automatically superior to parallel batch experimentation.
Hybrid platforms can use parallel preparation with serial reaction or analysis when that division matches the bottleneck.

\section{High-throughput kinetics: Mayer}

Kuffer and Mayer combined liquid-handling robotics with continuous UV monitoring to measure acid- and base-catalysed carbonyl-addition kinetics across a multiparameter space.\mcite{KufferMayer2026}
The workflow generates complete kinetic traces and converts them into rate constants rather than recording only a final absorbance.
A mechanistic rate model then separates specific acid, specific base, general acid or base, and background contributions.

The relative reactivity of the aldehydes changes between catalytic regimes.
That observation would be obscured by a dataset that treated every measurement as one undifferentiated prediction problem.
The case shows how throughput and mechanistic structure can reinforce one another when the measurement retains the relevant information.

\section{Automated synthesis and analysis: Bourne}

Avila and co-workers integrated a stopped-flow reactor into a high-throughput platform for library synthesis and at-line analysis.\mcite{Avila2022}
Reaction segments are moved into the reactor, held for a defined reaction time, and then transferred for analysis.
The architecture retains access to elevated temperature and pressure while reducing the material required to maintain a continuously operating steady state.

The campaign performed hundreds of reactions and generated a dataset for a feed-forward neural network.
The modelling component is useful because it consumes standardised observations produced by the platform.
The main Chapter~7 lesson is that reactor control, sample handling, analysis, and data capture must be engineered together before a model can guide later experiments reliably.

\section{Cage discovery: Greenaway progression}

\subsection{Computation fused with robotic synthesis}

Greenaway and co-workers combined computational screening with robotic synthesis for 78 precursor combinations targeting organic cages.\mcite{Greenaway2018}
The complete matrix records clean cage formation, partial formation, and negative outcomes rather than displaying only isolated successes.
Computational topology and formation-energy predictions can therefore be compared with the experimental result for every planned combination.

The screen produced clean formation of 33 cages and exposed unexpected cage-catenane chemistry.
Systematic execution did not remove surprise.
It made the surprise visible against a complete and internally comparable background.

\subsection{Automated synthesis, characterisation, and analysis}

Basford and co-workers expanded the workflow to 366 imine-condensation reactions and automated substantial parts of synthesis, sample preparation, characterisation, and data analysis.\mcite{Basford2024}
Turbidity, NMR, and high-resolution mass spectrometry report different aspects of the reaction outcome.
Their combination supports assignments of species type, conversion, and cage topology.

The study also links experimental results to a much larger set of computed hypothetical structures and pore properties.
Experiment maps, analysis code, and raw data were deposited alongside the article.
This combination makes the case especially useful for discussing practical FAIR data, provenance, and the importance of reusable analysis rather than throughput alone.

\section{Szczypiński cage automation}

Shields and co-workers combined low-cost conformational calculations with a semi-automated high-throughput crystallisation campaign for an amide-linked organic cage.\mcite{ShieldsCage2024}
The calculations identified several accessible cage conformations with different internal acid--acid distances.
The crystallisation screen then trapped five of the predicted conformers in experimentally determined solid-state structures.

This is a genuine automation case because a programmable platform executes a broad crystallisation campaign reproducibly.
It is not an autonomous case because experimental measurements do not select subsequent crystallisation conditions during the demonstrated workflow.
The distinction prepares the transition to the final chapter.

\section{From execution to autonomy}

Chapter~7 has assembled the infrastructure required to execute experiments, measure outcomes, perform quality control, and create machine-readable records.
Those capabilities are necessary but not sufficient for autonomy.
A loop becomes autonomous when the measurement changes the decision about which action or experiment occurs next.

\section{Summary}

Experiment count and information content are not equivalent.
FAIR chemical data require identities, conditions, measurement context, processing, provenance, and explicit failure states.
High-throughput batch and flow systems should be selected according to the chemical and analytical task.
The Mayer, Bourne, Greenaway, and Szczypiński cases show different ways to generate dense, reusable, and chemically informative data.
Automation becomes autonomy only when measured outcomes influence subsequent decisions.
